\begin{abstract}
A run-to-run standard error for a benchmark average built from per-benchmark variances alone ignores the covariance between benchmark scores. Our estimator SNAP splits each benchmark's items into two halves and takes cross-half products of run-centred scores, which removes item-sampling noise the halves don't share, although an item effect shared across benchmarks and halves would enter the estimate, and a 0.124 share of item variance would reproduce our margin value. We apply SNAP to 375 released DataDecide runs at the largest checkpoint step shared within each of 125 configurations, whose replicates differ in seed and partly in checkpoint schedule. On per-byte margins over ten benchmarks, the average's standard deviation is 1.244 times its independent value with interval \ci{1.143}{1.338}, and BoolQ alone carries 68\% of the covariance trace. Our accuracy interval \ci{0.993}{1.157} includes one under a low-powered design. A registered test on four unseen tasks fails its lower-limit rule at 1.217 with interval \ci{0.872}{1.498}. \outcome{\Ronecase}{1}{R1 pass}{\outcome{\Rtwocase}{1}{R2 not supported}{Under rules fixed before scoring, nine seeds at five PolyPythias sizes give margin inflation of \pending{R1 lambda} with interval \ci{\pending{lo}}{\pending{hi}}, and pairing that battery with a disjoint item bank leaves \pending{cross lambda} with the shared-item explanation rejected.}\outcome{\Rtwocase}{2}{R2 supported}{Under a rule fixed before scoring, nine seeds at five PolyPythias sizes give margin inflation of \pending{R1 lambda} with interval \ci{\pending{lo}}{\pending{hi}}, although a disjoint item bank supports a shared item component.}\outcome{\Rtwocase}{3}{R2 undecided}{Under a rule fixed before scoring, nine seeds at five PolyPythias sizes give margin inflation of \pending{R1 lambda} with interval \ci{\pending{lo}}{\pending{hi}}, although a disjoint item bank leaves the source of that inflation undecided.}}\outcome{\Ronecase}{2}{R1 fail}{On a second family, nine seeds at five PolyPythias sizes give margin inflation of \pending{R1 lambda} with interval \ci{\pending{lo}}{\pending{hi}}, which fails a rule fixed before scoring at a design with 40 within-size contrasts, and a disjoint item bank \outcome{\Rtwocase}{1}{R2 not supported}{leaves \pending{cross lambda} with the shared-item explanation rejected}\outcome{\Rtwocase}{2}{R2 supported}{supports a shared item component}\outcome{\Rtwocase}{3}{R2 undecided}{leaves the source of that inflation undecided}.} Apart from the registered test and the two PolyPythias rules, every analysis is exploratory. We recommend paired replicate scores across the chosen battery with a benchmark-removal check.
\end{abstract}

\section{Introduction}
A benchmark average can change across training runs even when the evaluation items stay fixed. For an equally weighted average of $K$ benchmarks, independent run deviations give variance $K^{-2}\operatorname{tr}\Sig$, where $\Sig$ is their cross-benchmark covariance matrix. Covariance adds the off-diagonal terms to that variance, and assuming independence therefore weakens a comparison's error control. In simulation with margin-like covariance, such a comparison declares a difference in 0.113 of replicates with a true gap of zero, against a nominal rate of 0.05. On DataDecide's recipe comparisons, whose runs within a size band share seed labels, the paired-difference standard deviation has median ratio 1.024 to its independence value and the wrong-call rate moves by at most 0.009, which means the covariance matters more for one battery average's error bar than for seed-matched comparisons.

We estimate item-general covariance from two disjoint item halves per benchmark, centring both half scores across the replicate runs within each configuration and taking their cross-half products. Under zero-mean item noise and conditional independence across halves, they estimate run covariance, including its diagonal. We adapt the repeated-measurement construction of quantitative genetics to this setting and call it SNAP, Seed Noise Across Phenotypes (Appendix~\ref{app:derivations}).

Our empirical study uses DataDecide \citep{magnusson2025}, with 25 data recipes and five model sizes that each have three released replicates, and those 125 configurations supply 375 runs across ten benchmarks, scored at the largest checkpoint step each configuration's replicates share. Our primary analysis reads released item scores without model inference and compares per-byte log-likelihood margins with accuracy on the same items. These two scales respond differently to changes in confidence, and margins carry more seed signal, with per-benchmark reliability averaging 0.652 against 0.324 for accuracy (Section~\ref{sec:composition}), which is why we treat margins as the primary scale and read accuracy as a bound.

Margin inflation $\Lam$, the ratio of the average's standard deviation to its independent value, is 1.244 with recipe-cluster interval \ci{1.143}{1.338}, which excludes one, while accuracy gives 1.078 with interval \ci{0.993}{1.157}, which includes one at the scored checkpoint (Table~\ref{tab:primary}).

The accuracy interval no longer includes one under five of the 25 recipe removals, removal of the 750M band, and removal of WinoGrande or BoolQ, while a simulation at the fitted accuracy cell would detect a true ratio of 1.10 in only 0.135 of replicates.

Because we built our estimator on these ten benchmarks, we registered a test on SciQ, MedMCQA, and multiple-choice DROP and CoQA on the same checkpoints at 530M, 750M, and 1B, whose rule requires the lower interval limit for held-out margin inflation to exceed one. The test fails on all 225 runs, where the held-out estimate is 1.217 with interval \ci{0.872}{1.498}. The original ten benchmarks on the same runs give 1.240 with interval \ci{1.054}{1.397} and pass it, a comparison fixed before any held-out score existed (Section~\ref{sec:heldout}).

A single model family and a seed label that also changes the training batch leave open whether the inflation belongs to DataDecide's design. We therefore scored all 45 PolyPythias runs at five sizes on the same 37,682 items, where each of nine seeds reruns one recipe, and paired that battery with a disjoint bank of 6,808 items built at the same OLMES commit. Two rules fixed before scoring read the results, and Section~\ref{sec:family} reports that margin inflation there is \pending{R1 lambda} with interval \ci{\pending{lo}}{\pending{hi}} and that the shared-item explanation is \outcome{\Rtwocase}{1}{R2 not supported}{rejected}\outcome{\Rtwocase}{2}{R2 supported}{supported}\outcome{\Rtwocase}{3}{R2 undecided}{undecided}.

This paper makes four contributions to measuring run covariance. We define the item-general run covariance that cross-half products estimate, identified up to an item component shared across benchmarks and halves, and state which assumptions the released runs can't be shown to satisfy. Algorithm~\ref{alg:snap} gives the measurement recipe with its interval calibration, where wild recipe-cluster coverage runs from 0.928 to 0.954 across twelve simulated populations. We measure it on two score scales across the 375 released runs, where margin inflation is 1.244 at a magnitude that BoolQ sets, and every recipe, recipe-pair and benchmark deletion, along with four of the five size-band deletions, leaves an interval that excludes one (dropping 530M gives a lower limit of 0.993). At that value the ten benchmarks average like 6.5 independent ones of the same average variance (interval 5.6 to 7.7), a within-battery count that doesn't rank batteries, since removing BoolQ raises the factor to 1.786 while the aggregate standard deviation falls from 0.00572 to 0.00516.

We register a test on unseen tasks that fails at its registered scope, where the held-out margin interval runs 0.626 wide against 0.343 for the original ten benchmarks on the same runs. In practice, a correction estimated on the same battery returns the simulated false-positive rate from 0.113 to 0.051 and in a separate transfer simulation leaves 0.101 at a battery whose own ratio is 1.5. Apart from the registered test and the two PolyPythias rules, every analysis is exploratory.
